\subsection*{Retrieving word images from earlier answers}
A visual writer can preserve a familiar word in memory and still select the wrong word for a new sentence. We added 2,048 fitting-answer word images to the masked writer and tested it on 147 previously unused conversational passages (Fig.~\ref{fig:word-memory}). On 2,238 questions the frozen parent produced 500 exact answers, a revised decoder with retrieval disabled produced 502, and enabling retrieval also produced 502. The last count fell to 59 after exchanging the question and history, and to 254 after exchanging the passage. The corresponding byte writer produced 473 exact answers with retrieval enabled.

None of the three visual conditions matched a complete reference on the 259 questions requiring new multiword phrasing. Visual word F1 with retrieval was 14.79\%, compared with 14.40\% for the original writer (paired difference 0.39 percentage points, 95\% story-bootstrap interval $-0.42$ to 1.28). Given the answer length and every other reference word, the writer recovered 356 of 1,100 available source words and none of 208 source-absent words. The correct image was stored for 108 of those 208. Among the 2,048 memory entries its median rank was 117, it ranked first three times, and it reached the first eight on 17 occasions. The median probability of using the retrieval route on those 108 slots was 0.106. Forcing retrieval would therefore recover only three of these missing words even with supplied neighbours.

The model wrote \emph{University of Bristol} when asked where Mr.\ Preece worked. Its answer changed under both input exchanges. For \emph{What was their deal?}, it wrote \emph{he had with palms}, against \emph{would toss a coin}. Retrieval also changed one answer from \emph{No} to the correct \emph{yes}, but that \emph{yes} persisted after exchanging the question. The measured failure concerns contextual selection among writable words. These aligned English print tests leave historical-script transfer and a compute advantage at matched language quality open.
